\chapter{Software Architecture and Design}\label{chap:design}

This chapter documents the structure of SHEO as a software design description in the
sense of IEEE~1016~\cite{ieee1016}. Where the
requirements chapter states \emph{what} the AI-Driven Smart Home Energy Optimizer
(SHEO) must do, this chapter states \emph{how} the system is structured to deliver
it: the design goals that govern every decision (\S\ref{sec:design-goals}), the
architectural style and its alternatives (\S\ref{sec:architecture}), the component
and class decomposition (\S\ref{sec:components}--\S\ref{sec:classes}), the
design patterns that give the structure its flexibility
(\S\ref{sec:patterns}), the detailed behaviour of the two central control paths
(\S\ref{sec:behaviour}), and finally the two subsystems that constitute the
intellectual heart of the product (\S\ref{sec:heart}). Throughout, the design
transforms the analysis model of the requirements chapter into a concrete design model.

\section{Design goals}
\label{sec:design-goals}

The design is governed by four fundamental design concepts: abstraction,
modularity, the cohesion/coupling balance, and information hiding. These are not
abstract ideals here; each is realised by a concrete mechanism in SHEO. The system
is decomposed into modules that each own a single capability (high cohesion) and
communicate through narrow, intention-revealing interfaces (low coupling), so that
a change to one capability does not ripple across the others. Abstraction is applied
twice over: device behaviour is abstracted behind a uniform adapter interface, and
device \emph{description} is abstracted into a declarative capability schema, so that
neither the user interface nor the validation logic ever names a concrete device
model. Information hiding is the organising principle of the data tier: every detail
of persistence is sealed behind repository interfaces, so that business logic remains
ignorant of, and therefore independent of, the storage technology.

Table~\ref{tab:design-goals} maps each design concept to the SHEO mechanism that
embodies it. This mapping is what makes the qualities verifiable rather than merely
asserted.

\begin{table}[htbp]
\centering
\caption{Each design concept and the SHEO mechanism that realises it.}
\label{tab:design-goals}
\begin{tabularx}{\textwidth}{L{0.20\textwidth} Y Y}
\toprule
\textbf{Design concept} & \textbf{What it requires} & \textbf{SHEO mechanism} \\
\midrule
Abstraction &
Suppress per-instance detail behind a general interface. &
The device-adapter interface abstracts device behaviour; the capability
schema abstracts a device type into data, so callers never branch on a concrete
model. \\
Modularity &
Partition the system into independently buildable, testable parts. &
A layered backend (API $\to$ services $\to$ repositories $\to$ domain) plus one
cohesive service per capability (devices, rules, optimisation, recommendations,
telemetry). \\
High cohesion / low coupling &
Each module has one reason to change; modules interact narrowly. &
Each service owns exactly one capability; producers and consumers are decoupled by
an in-process event bus, so a new consumer is added without touching producers. \\
Information hiding &
Hide volatile design decisions behind stable interfaces. &
Repositories hide all query construction and the storage engine; the capability
schema hides device-type variance from both the front-end and the capability-validation routine. \\
\bottomrule
\end{tabularx}
\end{table}

\section{Architectural design}
\label{sec:architecture}

\subsection{The architectural style}

SHEO is a \textbf{three-tier client--server} application whose server interior
composes three further architectural styles, each selected to satisfy
a specific quality goal:

\begin{enumerate}[leftmargin=1.4em]
  \item \textbf{Three-tier client--server.} Presentation (a browser-based
  single-page application), application logic (a Python application server), and data
  (a relational store) are cleanly separated, giving independent deployability and a
  single, well-defined transport boundary over HTTPS and secure WebSockets.
  \item \textbf{Layered architecture} inside the server. Dependencies point
  \emph{strictly downward}: API routers depend on services, services on
  repositories, repositories on the domain model and the object-relational mapping.
  No router issues a database query and no service imports the web framework, so each
  layer is replaceable and unit-testable in isolation.
  \item \textbf{Repository (data-centred) style} at the persistence boundary. Query
  construction is hidden behind intention-revealing methods, which also enforce
  unit-scoping for access control; a tariff repository additionally falls back to
  the building-wide tariff when a unit has no override.
  \item \textbf{Event-driven (Observer) integration} inside the server. An in-process
  event bus decouples event producers (the simulator, the rule engine, the device
  service) from consumers (the live-feed broadcaster and the notification service),
  so new consumers can be added without modifying producers --- a direct application
  of the open/closed principle.
\end{enumerate}

The client itself follows an \textbf{MVC} flavour: the REST and WebSocket payloads
are the Model, the page components are the View, and the authentication context
together with the event handlers act as the Controller. The capability-driven control
component is a pure View rendered over the capability schema, which serves as its
Model. Figure~\ref{fig:arch} shows the resulting three-tier and layered structure,
and the strict downward dependency rule that governs it.

\fig{fig02-architecture.png}{Three-tier and layered architecture. Dependencies point
strictly downward (presentation $\to$ application $\to$ data); no router runs a
database query and no service imports the web framework.}{fig:arch}

\subsection{The strict downward-dependency rule}

The single invariant that keeps the layered design honest is that a layer may depend
only on the layers beneath it, never on a peer or a layer above. The presentation
tier knows the application tier but not the database; the application tier knows the
repository abstraction but not the SQL dialect behind it; the data tier knows nothing
of the web. This invariant is what makes the headline maintainability claim true: a
change of storage engine is localised to the data tier, and a change of transport is
localised to the presentation tier, because no business rule has been allowed to leak
into either.

\subsection{Alternatives considered}

Two architectural decisions deserve an explicit record of the alternatives weighed.

\begin{itemize}[leftmargin=1.4em]
  \item \emph{In-process event bus versus an external message broker.} A broker
  (for example a message queue or an MQTT bus) would permit horizontal scaling but
  would force an external infrastructure dependency that a single-node deployment
  cannot assume. SHEO instead runs its simulator and scheduler as in-process
  asynchronous tasks that communicate over an in-process publish/subscribe bus.
  Because producers and consumers already talk through that publish/subscribe seam,
  a future move to a real broker would replace only the transport, not the
  producers or consumers.
  \item \emph{Layered monolith versus microservices.} Microservices were rejected as
  over-engineering for the scope. A layered monolith delivers most of the modularity
  benefit --- clear seams, independent testability, localised change --- at a small
  fraction of the operational cost, which is the correct trade-off on the
  function/cost/time triangle for a system of this size.
\end{itemize}

\subsection{Deployment}

The demonstration deployment is a \textbf{single application node} hosting the
application server, together with a single embedded database file on a local volume;
the browser is the only client node. The simulator and the scheduler are not separate
processes but background tasks within the application's lifespan --- which is exactly
why the simulator can publish telemetry straight onto the same in-process bus that the
WebSocket broadcaster listens on, with no broker in between. A production variant would
swap the embedded database for a networked relational server and run several worker
processes behind a TLS-terminating reverse proxy; the layered design localises that
change to the data tier and the deployment topology, leaving the application logic
untouched. Figure~\ref{fig:deploy} shows this topology.

\fig{fig03-deployment.png}{Deployment diagram. A single application node serves the
browser-based front-end over HTTPS and secure WebSockets and persists to a local
embedded database; the simulator and scheduler are in-process background tasks, not
separate hosts.}{fig:deploy}

\section{Component decomposition}
\label{sec:components}

The decomposition opens the black box of the use-case view and divides the system
along two orthogonal axes: \emph{by tier/layer} (presentation, application, data),
which yields the strict downward dependency flow, and \emph{by cohesion}, so that
each application-tier service owns exactly one capability and therefore has a single
reason to change. Figure~\ref{fig:package} shows the resulting package structure,
listed in dependency order so that each package depends only on those below it.

\fig{fig05-package.png}{Package and layered decomposition. Each package may depend
only on those below it --- the strict downward flow of the layered
architecture.}{fig:package}

The principal packages are as follows. The \textbf{presentation layer} holds the thin
routers and the dependency-injection helpers; each router validates a request,
resolves the authenticated user, calls exactly one service method, and returns a
data-transfer object. A small set of Administrator-only, read-only routes
(\texttt{GET /api/admin/overview} and the per-unit \texttt{/api/admin/units/\{id\}}
drill-in routes) sits alongside the per-unit routers and is guarded by the role
dependency. The \textbf{services layer} holds all business logic, one
cohesive service per capability, including the Administrator service that assembles
the building overview by reusing the per-unit telemetry service. The \textbf{cross-cutting layer} holds cross-cutting
concerns: the event bus, the security primitives, a single clock source, and a
framework-free domain-error hierarchy. The \textbf{device-adapter package} holds the
device-adapter interface hierarchy that encapsulates per-device-type behaviour. The
\textbf{simulator} and the \textbf{scheduler} are the two background workers --- the
former generating telemetry, the latter firing rules and accruing savings. The
\textbf{persistence layer} is the only layer that issues database queries, and the
\textbf{domain layer} holds the capability schema, the enumerations, and the
object-relational aggregates that form the structural foundation.

Crucially, the front-end reaches the server through exactly \textbf{two ports} and no
others: a \textbf{REST} port for request/response interaction, and a
\textbf{WebSocket} port for the live feed of telemetry and notifications. This narrow
contract is what allows the entire front-end to be developed and tested against a
stable, documented surface. Figure~\ref{fig:component} shows the component wiring:
routers depend on services through dependency injection, services depend on
repositories and adapters, and only the data tier touches the database. The event
bus is the integration hub --- the simulator and the device-control, rule, and
recommendation services publish to it, while the notification service and the
WebSocket connection manager subscribe.

\fig{fig07-component.png}{Component diagram. The front-end reaches the server only
through the REST and WebSocket ports. The event bus is the integration hub: the
simulator and the device-control, rule, and recommendation services publish, while
the notification service and the WebSocket connection manager subscribe.}{fig:component}

The component diagram above treats the front-end as a single box behind those two
ports. Figure~\ref{fig:frontend} opens that box and shows the front-end's own
internal structure, which realises the \textbf{MVC} flavour described in
\S\ref{sec:architecture}. The \emph{Controller} is \texttt{App.jsx} (routing and the
role gate) together with \texttt{auth/AuthContext.jsx} (the session, the bearer token,
and the current role); the \emph{View} is the navigation shell
(\texttt{Layout.jsx}), the page components, the capability-driven
\texttt{DeviceControl} component, and the shared UI primitives; and the
\emph{Model} is the data the client holds --- the REST and WebSocket payloads, the
capability schema that \texttt{DeviceControl} renders itself over, and the
\texttt{lib/format.js} helpers for currency, palette, and time. Crucially, the only
two arrows that leave the front-end are the same two ports as before:
\texttt{api/client.js} (REST) and \texttt{api/ws.js} (WebSocket). This keeps the
client's structure answerable to the same maintainability question as the server's
class diagram --- ``to change behaviour~X, which component do I touch?'' --- without
widening the contract to the server.

\fig{fig17-frontend-component.png}{Front-end component structure. The single-page
application follows an MVC split --- \texttt{App.jsx} and \texttt{AuthContext} as
Controller, \texttt{Layout}, the pages, \texttt{DeviceControl}, and the shared UI as
View, and the server payloads plus capability schema as Model --- and reaches the
server only through the same REST and WebSocket ports as
Figure~\ref{fig:component}.}{fig:frontend}

\section{Class design}
\label{sec:classes}

The static model, shown in Figure~\ref{fig:class}, is organised around the
\textbf{entity/boundary/control} separation. \emph{Entity} classes
(the object-relational aggregates, with the unit (a Home occupied by one Resident) as the aggregate root) hold
persistent state; \emph{boundary} classes (the routers and the capability-driven
control component) mediate interaction with the outside world; and \emph{control}
classes (the services) coordinate the work. This separation keeps presentation,
coordination, and state from contaminating one another.

\fig{fig06-class.png}{Class diagram. The unit (a Home occupied by one Resident) is the aggregate root; the
device-adapter interface hierarchy and the capability-schema value objects
drive device behaviour and command validation. The user--device control grant is
reified by an explicit permission class, and the rule-to-execution dependency is
deliberately non-cascading so that deleting a rule preserves its audit
trail.}{fig:class}

The control responsibilities are partitioned across six key services, each cohesive
and singly responsible:

\begin{itemize}[leftmargin=1.4em]
  \item \textbf{Device-control service} --- the single, uniform pipeline for device
  control and the \emph{single safety choke point}. Manual commands, rule-fired
  commands, and recommendation acceptances all flow through one command-application
  entry point, so the offline check, capability validation, and safety check cannot be
  bypassed by any caller. It returns a structured result indicating success, rejection,
  timeout, or a skipped no-op.
  \item \textbf{Rule engine} --- authoring, validation, conflict detection, evaluation,
  execution, and undo of the \textsc{when}--\textsc{then}--\textsc{until} automation
  rules. Validation checks an action against the device capability schema and detects
  time-overlapping opposing rules; execution is idempotent and rate-limited, dispatches
  through the device-control service only when auto-action is enabled, logs an immutable
  execution record, and supports a short undo window.
  \item \textbf{Optimisation/bill-saving service} --- the deterministic, explainable
  estimation engine. It computes the historical baseline, estimates a rule's monthly
  saving in Vietnamese đồng (VND), accrues measured savings into the current billing
  cycle, and flags rules whose measured saving drifts from the estimate.
  \item \textbf{Recommendation service} --- orchestrates the recommendation lifecycle:
  it delegates habit \emph{detection} to a swappable \textbf{recommendation-provider port}
  (the client's ``AI'', defaulting to a deterministic, explainable miner), then performs
  the parts that are SHEO's own software --- pricing each candidate in VND via the
  optimisation service, ranking and capping, turning an accepted suggestion into a rule,
  and suppressing a dismissed one by signature. A black-box model can replace the default
  provider behind the same port without affecting the VND estimator or any other layer.
  \item \textbf{Telemetry service} --- produces the dashboard aggregate, the
  unit-scoped consumption time-series, the top-consumers ranking, and the
  online/offline refresh that marks a device unreachable after a silence threshold.
  \item \textbf{Administrator service} --- serves the building owner's read-only oversight.
  It builds the \emph{building overview} (the resident roster with each unit's live power,
  cycle kWh, and bill, plus building-wide totals) and resolves a unit for the
  Administrator-only per-unit drill-in, raising a not-found error for an unknown unit. It
  reuses the per-unit telemetry service rather than duplicating its aggregation (DRY), and
  authors no rules and operates no devices.
\end{itemize}

\section{Design patterns}
\label{sec:patterns}

Five design patterns~\cite{gof} give the structure its flexibility; each is realised by a
concrete component and, where natural, advances a SOLID principle. Strategy, Observer,
and State are classical Gang-of-Four (GoF) behavioural patterns; Repository and
Dependency Injection are architectural patterns from Fowler's \emph{Patterns of
Enterprise Application Architecture} and the Dependency-Inversion Principle
respectively. Table~\ref{tab:patterns} records the pattern, its intent, and its
realisation in SHEO.

\begin{table}[htbp]
\centering
\caption{Design patterns realised in SHEO, with the SOLID principle each advances.
Strategy, Observer, and State are GoF behavioural patterns; Repository is an
architectural pattern; Dependency Injection follows the Dependency-Inversion
Principle.}
\label{tab:patterns}
\begin{tabularx}{\textwidth}{L{0.16\textwidth} L{0.30\textwidth} Y}
\toprule
\textbf{Pattern} & \textbf{Intent} & \textbf{Realisation in SHEO} \\
\midrule
Strategy (GoF) &
Encapsulate a family of interchangeable behaviours behind one interface. &
The device-adapter interface hierarchy (one strategy per device type) encapsulates each
type's command semantics and power model, so no code branches on type. Adding a device
type adds a strategy without editing callers --- the open/closed principle (OCP). The same
pattern isolates the ``AI'': habit detection sits behind a \emph{recommendation-provider}
port, so the default explainable miner and a future black-box ML provider are
interchangeable strategies, swapped at the composition root. \\
Observer (GoF) &
Notify dependents of state changes without coupling them to the source. &
The in-process event bus: producers publish typed events; the notification service and
the WebSocket broadcaster subscribe. New observers are added without touching
publishers. \\
Repository (Fowler) &
Mediate between the domain and the data source via a collection-like interface. &
A generic repository plus one repository per aggregate hides all query construction and
the storage engine, keeping business logic persistence-agnostic. \\
Dependency Injection (DIP) &
Supply a component's collaborators from outside rather than constructing them within. &
The dependency-injection helpers provide a ready database session, the current user, and
role guards to each router. Routers depend on abstractions rather than concrete
constructions --- a direct application of the Dependency-Inversion Principle, and the
structural basis of role-based access control. \\
State (GoF) &
Let an object alter its behaviour as its internal state changes. &
The rule and execution lifecycles are tracked by explicit status flags rather than
scattered ad-hoc conditions: a rule is enabled or disabled and may be flagged for
recalculation, and an execution is applied, undone, or past its undo window. Behaviour
then follows these recorded states. \\
\bottomrule
\end{tabularx}
\end{table}

\section{Detailed behaviour}
\label{sec:behaviour}

Two sequence diagrams capture the dynamic behaviour of the system's two principal
control paths. Figure~\ref{fig:seq-control} traces the \textbf{capability-driven
control} path. When a user operates a device, the front-end first fetches the
device's capability schema and renders the appropriate control widget from it; the
resulting command is sent to the server, where the device-control service authorises
the caller, validates the command against the same capability schema, runs the safety
check, asks the device's strategy adapter to apply the state change, publishes a
\emph{command-applied} event, and returns a structured result. The WebSocket
broadcaster, subscribed to the bus, then pushes the new state to every connected
front-end. The same path serves manual, rule-fired, and recommendation-driven
commands, which is what makes the safety guarantee non-bypassable.

\fig{fig10-seq-control.png}{Sequence: capability-driven device control. The command is
validated against the same capability schema used to render the control, then applied
through the device-adapter interface, and the new state is broadcast over the event
bus to all connected front-ends.}{fig:seq-control}

Figure~\ref{fig:seq-rule} traces the \textbf{automatic rule-firing} path driven by
the scheduler. On each tick the scheduler asks the rule engine to evaluate each
enabled rule; when a rule's \textsc{when} condition holds and the device is not
already in the target state, the engine dispatches the action through the
device-control service (only if auto-action is enabled for that rule), writes an
immutable execution record, and publishes a \emph{rule-fired} event. The
notification service, observing that event, persists a notification and re-publishes
it for the broadcaster, so the user sees the automated action and retains a short
window in which to undo it.

\fig{fig11-seq-rule-firing.png}{Sequence: automatic rule firing. The scheduler
evaluates an enabled rule; when its condition holds, the engine dispatches through the
same control pipeline, logs an immutable execution, and the firing is observed by the
notification service and broadcast to all connected front-ends.}{fig:seq-rule}

\section{The two heart subsystems}
\label{sec:heart}

Two subsystems carry the product's distinctive value and therefore warrant detailed
treatment: the capability schema, which makes the system extensible without
per-model code, and the bill-saving estimation algorithm, which makes every saving
figure transparent and verifiable.

\subsection{The capability schema: a single source of truth}

The capability schema is the design device that lets SHEO support an open-ended set of
device types without writing a screen or a validator for each. A device type is
described declaratively as data, not as code branches. A \emph{control} value object
describes one controllable feature by its name, label, and \emph{kind} --- one of
toggle, range, enumeration, or read-only --- together with the bounds (minimum,
maximum, step), the permitted values, the default, and a flag marking the control as
safety-sensitive. A \emph{capability schema} lists a device type's telemetry channels
and its controls, plus whether the type is reversible, whether it is safety-critical by
default, and its nominal power draw. A registry holds exactly one schema per device
type.

The decisive property is that this single description \textbf{drives two consumers at
once}. The front-end renders its control widgets directly from the schema --- a toggle
for a toggle control, a slider for a range, a segmented selector for an enumeration
--- so the user interface never names a concrete device model. The server validates
every incoming command against the \emph{same} schema before dispatching it, using
the same capability-validation routine on both the manual and the rule-fired path.
Because rendering and validation read one definition, they can never disagree: a
control the interface offers is exactly a control the server accepts, with exactly the
bounds the interface enforced. Adding a new device type is therefore a two-element
change --- one schema entry and one device adapter --- with no edit to any router,
service, or screen.

Table~\ref{tab:capability-example} illustrates the structure of a schema entry for a
representative device type (a room air-conditioner), showing the three controllable
features it exposes.

\begin{table}[htbp]
\centering
\caption{Capability schema entry for a room air-conditioner. Each row describes one
controllable feature: its kind determines how the front-end renders the widget; its
bounds constrain the accepted values; the safety flag causes the capability-validation
routine to apply an additional safety check before dispatch.}
\label{tab:capability-example}
\begin{tabularx}{\textwidth}{L{0.16\textwidth} L{0.15\textwidth} Y L{0.18\textwidth}}
\toprule
\textbf{Control} & \textbf{Kind} & \textbf{Bounds / permitted values} & \textbf{Safety-sensitive} \\
\midrule
Power & Toggle & \{on, off\}; default: off & No \\
Target temperature & Range & 16--30\,\textdegree C, step 1\,\textdegree C; default: 26 & Yes \\
Operating mode & Enumeration & \{cool, fan, dry, auto\}; default: cool & No \\
\bottomrule
\end{tabularx}
\end{table}

\subsection{The bill-saving estimation algorithm}

The second heart subsystem is the explainable estimation of a rule's monetary saving.
It is deliberately kept inside SHEO --- never behind the recommendation-provider's
black box --- so the figure the client values most stays transparent. Every estimate
reduces to one formula, summed over the hours a rule affects:

\begin{equation*}
\begin{split}
  \text{saved\_vnd}
  \;=\;
  \sum_{h \,\in\, \text{affected hours}}
  \bigl(&\,\text{baseline\_kWh}_h - \text{expected\_kWh\_with\_rule}_h\,\bigr) \\
  &\times \text{tariff\_VND\_per\_kWh}.
\end{split}
\end{equation*}

The estimate is computed in three transparent steps.

\begin{enumerate}[leftmargin=1.4em]
  \item \textbf{Baseline.} For the device in question, the system computes an
  \emph{hourly} profile of average energy use --- the mean kilowatt-hours consumed in
  each hour of the day --- over the previous \textbf{14 days}. Summing the profile over
  the hours the rule affects gives the baseline energy that the rule will act upon.
  \item \textbf{Expected use with the rule.} A transparent per-action model derives what
  the device would draw \emph{if the rule were in force}. Switching a device off yields
  zero expected energy; raising an air-conditioner target by a degree applies a fixed
  saving fraction per degree (capped); reducing a fan speed or a bulb brightness scales
  the energy by the ratio of the new setting to the current one. Each rule of thumb is
  documented, so a reader can reproduce the figure by hand.
  \item \textbf{Pricing.} The avoided energy --- the baseline minus the expected use,
  floored at zero --- is multiplied by the effective tariff price in VND per
  kilowatt-hour, resolved from the building-wide tariff set by the Administrator (flat,
  tiered, or time-of-use), falling back to it when a unit has no override. The daily
  saving is projected to a month for display.
\end{enumerate}

The design property that matters most is \emph{when} the figure is shown: the estimated
VND saving is presented to the user \textbf{before a rule is accepted}, so that a saving
is a promise the user can weigh, not a surprise discovered later. Each estimate is
accompanied by a plain-language explanation naming the baseline, the expected use, and
the price, closing the loop on the explainability constraint. After a rule is in force,
the scheduler accrues the \emph{measured} saving per billing cycle and flags any rule
whose measured saving drifts materially from its estimate, so the promise is
subsequently checked against reality.
